L'automazione delle operazioni colturali in orticoltura richiede sistemi di guida autonoma capaci di mantenere il veicolo centrato lungo l'interfilare con precisione centimetrica. Nei contesti orticoli (es. coltivazioni di lattuga, pomodoro e cetriolo, sia in pieno campo che in serra), i sistemi di posizionamento globale basati su GNSS e correzioni cinematiche RTK manifestano schermature del segnale satellitare e riflessioni multi-path provocate dalla densa volta fogliare o dalle armature metalliche delle serre, oltre a costi infrastrutturali significativi per stazioni base locali. I sensori LiDAR 3D forniscono misurazioni spaziali dirette ma presentano costi unitari elevati per l'adozione su flotte commerciali e producono nuvole di punti prive del contenuto cromatico e tessiturale necessario per distinguere il terreno calpestabile dalla vegetazione infestante.

Questa tesi presenta un'architettura di navigazione visuale reattiva basata su telecamera monoculare RGB a basso costo, operante in anello chiuso a 20~Hz su computer di bordo embedded (NVIDIA Jetson Orin Nano). L'architettura separa la generazione assistita del dataset (eseguita offline) dall'inferenza a bordo del robot in tempo reale. Per ridurre l'onere di etichettatura manuale, il foundation model Segment Anything Model (SAM~3) genera maschere binarie di corridoio interfilare guidate da prompt. A bordo del robot, le immagini acquisite a risoluzione $224 \times 224$ sono elaborate da una rete leggera MobileNetV3 con testa LR-ASPP esportata in ONNX Runtime FP16, con un tempo medio di inferenza inferiore a 15~ms.

Dalla maschera predetta, la traiettoria di guida è ricavata mediante un algoritmo di tracking geometrico bilaterale disaccoppiato. La regione d'interesse è suddivisa in 24 strisce orizzontali dal basso verso l'alto; i confini fisici sinistro e destro sono tracciati separatamente tramite gate di continuità spaziale e interpolati con minimi quadrati iterativamente ripesati (IRLS) e funzione di costo di Cauchy. La linea centrale è ricavata analiticamente per bisezione dei parametri convergenti.

Il metodo proposto è confrontato con tre baseline consolidati in letteratura: proiezione a istogramma di colonna (SegMin), regressione su centroidi a strisce e interpolazione di contorno tramite RANSAC. La tesi valuta i trade-off strutturali dei singoli approcci: il metodo a doppio bordo proposto elimina il jitter angolare stocastico tipico di RANSAC ($\sigma_\omega = 0.00^\circ$) e previene gli errori sistematici introdotti da intrusioni fogliari asimmetriche nel calcolo dei centroidi. L'analisi sperimentale identifica e formalizza inoltre i limiti intrinseci del metodo proposto: la dipendenza dal seed iniziale nelle strisce inferiori in presenza di occlusioni vicino al veicolo, la sensibilità delle soglie di salto al beccheggio e rollio su terreno dissestato, l'inadeguatezza del modello a retta su filari curvilinei o svolte di fine campo, e la perdita di traccia in caso di interruzioni prolungate della vegetazione.

La validazione ad anello chiuso è condotta sulla piattaforma robotica Demetra, un rover agricolo a quattro ruote motrici e sterzanti indipendenti (4WS). Il centraggio in campo è monitorato mediante un sensore laser ausiliario indipendente, confermando la fattibilità del controllo visuale deterministico a 20~Hz senza ricorrere a LiDAR o sensori satellitari.
